\scanpage{36}
% archive sheet number: 35
% L1–3
\begin{center}\rule{.88\linewidth}{.35pt}\\[-.85em]\rule{.88\linewidth}{.65pt}\end{center}
\textsc{Algèbre}. --- \emph{Définition des classes de Stiefel-Whitney d'un module quadratique sur un corps de caractéristique différente de $2$.} Note de M. \textbf{\textsc{Antoine Delzant}}, transmise par M. Jean Leray.

% L4–5
Soit $\mathrm A$ un corps de caractéristique différente de $2$. On ne considère que des modules quadratiques $(\mathrm E,\mathrm Q)$ $(^{1})$, de rang fini sur $\mathrm A$, et non dégénérés.

% L6–16
On désigne par $\mathrm Q^{+}$ l'ensemble des classes à isomorphisme près de modules quadratiques. L'ensemble $\mathrm Q^{+}$ est muni de deux opérations : l'une, la somme directe externe, définit sur $\mathrm Q^{+}$ une structure de monoïde abélien pour laquelle tous les éléments sont réguliers (théorème de Witt); l'autre, le produit tensoriel, est commutative, distributive par rapport à la précédente, et admet un élément unité [la classe du module $(\mathrm A,\mathrm I)$, où $\mathrm I(x)=x^{2}$]. L'ensemble $\mathrm Q^{+}$ se plonge donc injectivement dans un anneau $\mathrm Q(\mathrm A)$, qu'il engendre. L'anneau $\mathrm Q(\mathrm A)$, commutatif avec élément unité, est appelé l'\emph{anneau des classes de formes quadratiques}. A toute classe $q$ de module quadratique $(\mathrm E,\mathrm Q)$ on fait correspondre $\varepsilon(q)=rg\,(\mathrm E)$; on définit ainsi un homomorphisme d'anneau de $\mathrm Q(\mathrm A)$ dans $\mathbf Z$ appelé l'augmentation.

% L17–21
Soit $\mathrm A_{s}$ la clôture séparable de $\mathrm A$ et $\mathrm G$ le groupe de Galois de $\mathrm A_{s}$ sur $\mathrm A$; on note $\mathrm H^{q}(\mathrm A,\mathbf Z/2\,\mathbf Z)$ le $q^{\text{--ième}}$ groupe de cohomologie de $\mathrm G$ à valeurs dans $\mathbf Z/2\,\mathbf Z$. On note additivement le groupe $\mathrm A^{\star}$, groupe multiplicatif de $\mathrm A$, et l'on pose $\mathrm D(\mathrm A)=\mathrm A^{\star}/2\,\mathrm A^{\star}$. Le groupe $\mathrm D(\mathrm A)$ s'identifie à $\mathrm H^{1}(\mathrm A,\mathbf Z/2\,\mathbf Z)$.

% L22–27
Soit $a$ un élément de $\mathrm A^{\star}$; on note $e_{a}$ sa classe dans $\mathrm D(\mathrm A)$, et $\varphi_{a}$ la classe du module quadratique $[\mathrm A,\mathrm Q;\ \mathrm Q(x)=ax^{2}]$. L'élément $\varphi_{a}$ ne dépend que de $a$ modulo $2\,\mathrm A^{\star}$; on peut donc définir un homomorphisme de $\mathrm D(\mathrm A)$ dans le groupe multiplicatif $\mathrm Q(\mathrm A)^{\star}$ des éléments inversibles de $\mathrm Q(\mathrm A)$. D'où un homomorphisme d'anneaux :
\[
\varphi:\quad \mathbf Z(\mathrm D)\to\mathrm Q(\mathrm A),
\]
% L28–29
où $\mathbf Z(\mathrm D)$ est l'algèbre de $\mathrm D(\mathrm A)$ sur $\mathbf Z$, munie de la filtration définie par les puissances de l'idéal d'augmentation.

% L30–32
\textsc{Proposition}. --- \emph{L'application $\varphi$ est surjective; elle est compatible avec les augmentations; le noyau de $\varphi$ est l'idéal engendré par les éléments $e_{a}+e_{b}-e_{a'}-e_{b'}$, tels que $\varphi_{a}+\varphi_{b}=\varphi_{a'}+\varphi_{b'}$.}

% L33–35
Le cup-produit définit une application de $\mathrm D(\mathrm A)\times\mathrm D(\mathrm A)$ dans $\mathrm H^{2}(\mathrm A,\mathbf Z/2\,\mathbf Z)$ notée $(e_{a},e_{b})\to a.b$. Si l'on identifie $\mathrm H^{2}(\mathrm A,\mathbf Z/2\,\mathbf Z)$ au sous-groupe du groupe de Brauer de $\mathrm A$ formé des $x$ tels que $2x=0$, l'élément $a.b$
